\documentclass[10pt,a4paper]{amsart}
\usepackage[margin=19mm]{geometry}
\usepackage{amsmath,amssymb,mathtools}
\usepackage[scr=rsfs,cal=euler]{mathalfa}
\usepackage{xfrac}
\usepackage[unicode,hidelinks,bookmarks=false]{hyperref}
\newcommand{\ie}{i.e.\ }
\newcommand{\cf}{cf.\ }
\newcommand{\resp}{respectively\ }
\newcommand{\Subsection}{\subsection}
\providecommand{\nobreakdash}{\nobreak\hbox{-}}
\setlength{\emergencystretch}{2em}
\multlinegap0pt
\usepackage[shortlabels]{enumitem}
\usepackage{amssymb}
\usepackage{bm}
\usepackage{xcolor}
\newtheorem{lemma}{Lemma}
\newcommand{\N}{\mathbb N}
\newcommand{\dtau}{\,{\rm d} \tau}
\newcommand{\tstar}{t^*}
\title{On the stability of solutions to non-Newtonian Navier--Stokes--Fourier-like systems in the supercritical case}
\author{Anna Abbatiello, Miroslav Bul\'i\v{c}ek, Petr Kaplick\'y}
\date{Source: arXiv:2601.16868v1 (CC BY 4.0)}
\begin{document}
\maketitle

\noindent\emph{Source and licence.} On the stability of solutions to non-Newtonian Navier--Stokes--Fourier-like systems in the supercritical case. Version arXiv:2601.16868v1 (CC BY 4.0), distributed by arXiv under the Creative Commons Attribution 4.0 International licence, http://creativecommons.org/licenses/by/4.0/, as recorded in the arXiv licence metadata for this exact version and shown on the abstract page https://arxiv.org/abs/2601.16868v1. Full licence notice: \texttt{licenses/arxiv-2601.16868-CC-BY-4.0.txt}.\par\smallskip

\noindent\textit{Editorial scope.} Counted unit: the article's lemma lem:decay (source0.tex lines 2013--2037). The article's complete original proof of it is delivered with it (source0.tex lines 2038--2056, one \texttt{proof} environment of 19 lines). No mathematics is written, repaired or re-derived: the statement and the proof are the article's own lines, in the article's own notation, and no retained line is substituted or re-flowed.  No further source-local context is needed: every cross-reference of every retained line resolves inside the two delivered spans.  Nothing was omitted from any retained span, so the package carries no omission record.  No retained line contains a citation, so the wrapper carries no bibliography and no bibliography entry.  No retained line contains a figure environment, an image-inclusion command or drawing code, so no image is delivered and the package declares no asset dependency.  Editorial formatting only, in addition: an AMS \texttt{amsart} preamble that restates the article's own declarations and macros used by the retained lines, namely the environments \texttt{lemma} on separate counters, together with the standard packages those lines need.  The retained environments print in this document as Lemma 1 in order of appearance, whereas the article prints the same items with the numbering of its own full document.  The complete native source of the article as distributed in the arXiv e-print for this exact version is bundled unchanged as \texttt{sources/193271/source0.tex}.
\begin{lemma}[Decay estimates]\label{lem:decay} \ 
\begin{itemize}[left=0pt]
\item[1)]
Let $\tstar \geq 0$ and let $f : (\tstar, +\infty) \to [0, +\infty]$ be a measurable function such that, for almost all $s, t \in (\tstar, +\infty)$ with $s < t$, there holds
\begin{equation}\label{eq:0}
f(t) + C_1 \int_s^t f(\tau) \, \dtau \leq f(s).
\end{equation}
Then
$$
f(t) \leq f(s) \, e^{-C_1 (t-s)}
$$
for almost all $s, t > \tstar$ with $s < t$.

\item[2)] Let $\tstar > 0$, $\alpha \in (0,1)$, and let $f : (0, \tstar) \to [0, +\infty]$ be measurable such that, for almost all $s, t \in (0, \tstar)$ with $s < t$, there holds
\begin{equation}\label{eq:1}
\max\Bigl(1, f(t) + C_1 \int_s^t f(\tau)^\alpha \, \dtau \Bigr) \leq f(s).
\end{equation}
Then
$$
f(t) \leq \left(\frac{f(s)}{1 + C_1 (t-s)}\right)^{1/\alpha}
$$
for almost all $s, t \in (0, \tstar)$ with $s < t$.

\end{itemize}
\end{lemma}

\begin{proof}
  Ad 1)
  We redefine $f$ on a set of zero measure so that $f$ is decreasing on $(\tstar,+\infty)$. We denote the new function again $f$. It satisfies  \eqref{eq:0} for all $s,t>\tstar$, $s<t$. Since $f$ is decreasing we get from \eqref{eq:0} that 
  \begin{equation}\label{eq:2}
  f(\tau)\leq\frac{f(\sigma)}{1+C_1(\tau-\sigma)}
\end{equation}
for all $\sigma,\tau\in(\tstar,+\infty)$, $\sigma<\tau$.
Splitting interval $(s,t)$ to $n\in\N$ parts and applying \eqref{eq:2} iteratively gives
$$
  f(t)\leq f(s)\left(\frac1{1+C_1\frac{t-s}n}\right)^n.
$$
The claim follows by limit passage $n\to+\infty$ since we modified the function $f$ only on a set of zero measure.

Ad 2) As in 1) we redefine $f$ on a set of zero measure so that $f$ is decreasing and satisfies \eqref{eq:1} for all $s,t\in(0,\tstar)$. Since $f\geq 1$ on $(0,\tstar)$ we get for all $s,t\in (0,\tstar)$, $s<t$
  $$
  f(t)\leq \left(\frac{f(s)}{1+C_1(t-s)}\right)^{\frac1\alpha},
  $$
  which is the claim since $f$ was modified only on a zero measure set.  
\end{proof}

\end{document}
